\section{Discussion}
\label{sec:discussion}

\textbf{Verifiability as a measurable orchestration principle.} The routing and orchestration literature asks which model should serve which request and answers with learned cost-quality predictors or bandit feedback \cite{li2026progrouter,sun2026cborchestrator,koul2026register,cao2026taskspecific}; our results supply the ground truth such predictors ultimately need, and suggest that for scientific agents a large share of the answer is available \emph{before} any routing data is collected --- from the epistemic structure of the step itself. Within the annotated distribution, a step that produces an artifact with an objective correctness standard is predictable in its dependence, and a step whose quality cannot be objectively adjudicated is comparatively safe to delegate; our held-out and true cross-domain checks (Results) show this regularity does not extend beyond the measured distribution --- it reverses on unseen instances of the training tasks and is consistently negative on an unseen domain --- so the principle currently licenses annotation-guided allocation, not untrained automatic routing. This reframing matters for practice: verification-oriented scaffolds and test-generation agents are already being built \cite{lee2025learning,kapetanovic2026phased}, domain-specific agent stacks for materials discovery expose the same frontier-dependence question in another guise \cite{rothfarb2025hierarchical}, component-level ablations of data-science agent workflows show how much of end-to-end performance is attributable to individual modules \cite{huang2026llmgenerated}, and surveys of autonomous discovery systems converge on orchestration as the open problem \cite{zhou2025autonomous}; our measurement protocol provides the missing evaluation layer for all of these --- not "does the small model feel adequate" but "does the small model, under identical context, reproduce frontier outcomes at rates that matter". The finding that naive splicing amplifies damage adds a cautionary boundary: step boundaries in current agent frameworks are not clean interfaces, and hybrid frontier/cheap pipelines must be engineered around recovery information, not assumed to compose.

\textbf{Relation to process supervision.} Process reward models and step-level verifiers pursue the same granularity by \emph{learning} step scores from data \cite{zhang2025groundedprm,zhang2025linking,rahman2025scaleprm,samarinas2026truncated,liu2026confspec,zhou2026silr}, with recent work documenting both reward hacking and faithfulness failures of learned verifiers \cite{ma2024what,nie2026stop,zhang2026when}. Our protocol occupies the complementary extreme: in an environment where success is objectively executable, step-level ground truth can be \emph{measured} rather than learned, at the price of executing the counterfactual. The two approaches are not competitors but stages of a pipeline --- measured replacement outcomes of the kind we produce are exactly the training signal that learned process models lack, and our blind-coded features are exactly the cheap proxies such models could be distilled from --- complements to rubric-based reward ranking and dense verifiable-reward designs that enrich the reward signal without a learned verifier \cite{li2026rlvr,wang2026beyond}. The idiosyncratic error family we observe (a frontier-quality open-weights model failing signal steps in ways neither API model does) also carries a message for the verifier literature: step failures are not draws from a shared difficulty distribution, and verifiers trained on one model's failures may be structurally blind to another's --- a heterogeneity that aggregate pass-rate evaluation cannot surface. Our held-out and cross-domain checks add the complementary warning: blind-coded features, like learned verifiers, fail to transfer outside their sampling regime --- the mechanism is not statistical noise but a reversed feature--outcome mapping for deliverable steps in the new domain --- so neither route yields a free out-of-distribution step scorer, and each new domain requires its own measured ground truth.

\textbf{Limits and what would falsify the picture.} Our evidence comes from one agent skeleton, one task domain (scientific computing on tabular data), and 74 pre-registered step instances (49 in-distribution, 25 out-of-distribution); per-cell statistical power is modest enough that no individual step-arm contrast survives multiple-comparison correction, and the headline claims rest on aggregations (difference-in-differences, error-family counts, cross-validated prediction) rather than single contrasts. The in-distribution prediction result now meets both pre-registered criteria at the enlarged sample size (49 points), and the out-of-distribution failure is established at triple the initial held-out power, so neither direction any longer depends on small-sample caveats. The verifiability features are blind-coded but human-defined, one (artifact production) is an operational proxy rather than a semantic property, one annotator is an AI assistant whose contamination we bounded but could not eliminate, and both out-of-distribution tests show the cross-validated gain does not transfer --- to unseen instances of the training tasks at either the instance or the task level, nor to a fully unseen domain where the outcome mapping for deliverable steps reverses. The strongest open test is no longer generality of the prediction but closing the loop from prediction to live orchestration decisions, and scaling the measured ground truth across domains --- each new domain costs one annotation cycle under the frozen protocol, a budget the cross-domain experiments show to be well spent. Within these bounds, the picture is coherent and actionable: frontier capability in scientific agents is a step-level, mechanistically heterogeneous, and partially predictable quantity, and the tools to measure it cost roughly a hundred dollars of compute --- orders of magnitude below the cost of the deployments whose design they should inform. Evaluating agents at that resolution is, we argue, the cheapest available upgrade to the science of agent engineering, and a necessary one if the field is to keep its claims about "where intelligence is needed" honest \cite{beel2025evaluating,miyai2025scientist,emami2026position}.

